\documentclass[12pt]{article}
\usepackage[T1]{fontenc}
\usepackage[utf8]{inputenc}
\usepackage{lmodern}
\usepackage{textcomp}
\usepackage{enumitem}
\usepackage{geometry}
\geometry{margin=1in}
\begin{document}
\begin{center}
Answer Key - History - Undergraduate Level Assessment 14 \
\small (Answers are ordered exactly as in the question paper.)
\end{center}

\section*{Multiple Choice}
\begin{enumerate}[label=\arabic*.]
\item \textbf{Answer:} C
\\ \textit{Options:} A) Iberomaurusian culture; Maglemosian; B) Backed blade; microlith blade; C) Iberomaurusian culture; Capsian culture; D) Natufian; Iberomaurusian culture; E) Capsian culture; microlith blade
\\ \textit{Note:} The correct answer reflects the cultural transition in Holocene Africa where the Iberomaurusian culture, characterized by specific lithic technologies and subsistence strategies, was succeeded by the Capsian culture, which introduced new advancements in tool-making and a shift in dietary practices.
\item \textbf{Answer:} A
\\ \textit{Options:} A) advanced metallurgy.; B) monumental works.; C) a system of money.; D) maritime navigation technology.; E) complex religious systems.
\\ \textit{Note:} The correct answer is advanced metallurgy because ancient South American civilizations, such as the Incas and Moche, primarily utilized gold, silver, and copper for decorative and ceremonial purposes rather than developing extensive metallurgical techniques comparable to those of contemporary civilizations in other regions.
\item \textbf{Answer:} A
\\ \textit{Options:} A) largest; B) most populated; C) earliest; D) least known; E) richest
\\ \textit{Note:} Tiwanaku was one of the largest and most influential urban centers of the ancient Andean civilization, known for its extensive architecture, agricultural innovations, and vast trade networks, which contributed to its significant size and power in the region.
\item \textbf{Answer:} D
\\ \textit{Options:} A) meat to fruits and vegetables.; B) species focused on grasses to fruits and vegetables.; C) nuts and fruits to fish.; D) nuts and fruits to species more focused on grasses.; E) species focused on grasses to meat.
\\ \textit{Note:} Research indicates that the evolution of the C$_{4}$ photosynthetic pathway in grasses allowed for more efficient growth in warmer, drier climates, leading ancient hominids to shift their diets from the more easily gathered nuts and fruits to a greater reliance on grass-based resources, which were becoming more abundant.
\item \textbf{Answer:} C
\\ \textit{Options:} A) Mayan; B) Kwakiutl; C) Mississippian; D) Inca; E) Clovis
\\ \textit{Note:} The correct answer is the Mississippian because Hernando de Soto's expedition in the 1540 s encountered complex, agricultural societies in the southeastern United States that were descendants of the Mississippian culture, known for their mound-building and advanced social structures.
\end{enumerate}

\section*{True / False}
\begin{enumerate}[label=\arabic*.]
\setcounter{enumi}{5}
\item \textbf{Answer:} True
\\ \textit{Options:} A) True; B) False
\\ \textit{Note:} According to the passage: Megalopolis
\item \textbf{Answer:} False
\\ \textit{Options:} A) True; B) False
\\ \textit{Note:} The correct answer is: Ottonian
\item \textbf{Answer:} True
\\ \textit{Options:} A) True; B) False
\\ \textit{Note:} According to the passage: 12.5
\item \textbf{Answer:} False
\\ \textit{Options:} A) True; B) False
\\ \textit{Note:} The correct answer is: 88
\item \textbf{Answer:} True
\\ \textit{Options:} A) True; B) False
\\ \textit{Note:} According to the passage: the proper aim point automatically
\end{enumerate}

\section*{Long Form Response}
\begin{enumerate}[label=\arabic*.]
\setcounter{enumi}{10}
\item \textbf{Answer:} Land for peace
\\ \textit{Note:} Expected answer: Land for peace
\item \textbf{Answer:} Journalist and author Thomas Laird
\\ \textit{Note:} Expected answer: Journalist and author Thomas Laird
\end{enumerate}

\vfill
\noindent\textit{End of Answer Key}
\end{document}